\documentclass[12pt]{article}

\usepackage[a4paper, margin=1in]{geometry}
\usepackage{xcolor}
\usepackage{tcolorbox}
\usepackage{fontawesome5} % icons
\usepackage{hyperref}
\usepackage{graphicx}
\usepackage{comment}
\usepackage[nameinlink,capitalize,noabbrev]{cleveref}
\usepackage{subcaption} % modern replacement for the deprecated subfigure package

\setlength{\parindent}{0pt}

% --- COLOR PALETTE (customize these) ---
\definecolor{maincolor}{HTML}{2A9D8F} % your teal
\definecolor{accentcolor}{HTML}{E76F51}
\definecolor{lightcolor}{HTML}{FFDEAD}

% Tell cleveref how to name your counter
\crefname{exercise}{assignment}{assignments}
\Crefname{exercise}{Assignment}{Assignments}

% --- ASSIGNMENT STYLE ---
\newcounter{exercise}

% Optional label: \exercisetitle[<label>]{<Title>}
\newcommand{\exercisetitle}[2][]{%
  \refstepcounter{exercise}% make this step referenceable/hyperlinked
  \vspace{0.5cm}%
  \colorbox{maincolor}{%
    \makebox[\dimexpr\linewidth-2\fboxsep][l]{%
      \textcolor{white}{\faPen\ Assignment \theexercise:~ #2}%
    }%
  }%
  % If an optional label was provided, set it
  \if\relax\detokenize{#1}\relax\else\label{#1}\fi
  \vspace{0.3cm}%
}


% --- IMPORTANT BOX ---
\newtcolorbox{importantbox}{
  colback=lightcolor,
  colframe=accentcolor,
  arc=4mm,
  boxrule=1pt,
  left=6pt,
  right=6pt,
  top=6pt,
  bottom=6pt
}

% --- PROMPT BOX ---
\newtcolorbox{promptbox}{
  colback=white,
  colframe=maincolor,
  arc=4mm,
  boxrule=1pt,
  left=6pt,
  right=6pt,
  top=6pt,
  bottom=6pt
}

% --- QUESTION BOX (with starred variant and cleveref names) ---
% Requires: \usepackage{hyperref} for \Form/\TextField
%           \usepackage{fontawesome5} for \faEdit

% Cleveref names for the question counter
\crefname{question}{question}{questions}
\Crefname{question}{Question}{Questions}

\makeatletter
\newcounter{question}
\newlength\questionheight
\setlength\questionheight{5cm} % default answer box height

% \question{<text>}                  -> with TextField (default height)
% \question[<height>]{<text>}        -> with TextField (custom height)
% \question*{<text>}                 -> no TextField, still numbered
\newcommand{\question}{\@ifstar{\question@text}{\question@field}}

% Common heading (increments counter and creates ref anchor)
\newcommand{\question@heading}[1]{%
  \refstepcounter{question}%
  \vspace{0.3cm}%
  \noindent\textbf{\faEdit\ Question \thequestion: #1}%\\[0.3em]%
}

% Starred form: numbered question without input field
\newcommand{\question@text}[1]{%
  \question@heading{#1}%
%  \vspace{0.3cm}%
}

% Unstarred form: numbered question with input field (optional height)
\newcommand{\question@field}[2][\questionheight]{%
  \question@heading{#2}%
 \\[0.3em] \begin{Form}
    \TextField[
      name=Q\thequestion,
      multiline=true,
      width=\linewidth,
      height=#1,
      charsize=10pt,
      bordercolor={0 0 0},
      borderwidth=1.8,
      backgroundcolor={1 1 1}
    ]{}
  \end{Form}
  \vspace{0.3cm}%
}
\makeatother

% --- SOLUTION BOX ---
%--- Definition of solution/no solution env to change between the teacher and the student versions --------%

% --- Teacher toggle ---
\newif\ifteacher
\teachertrue  % Write teacherfalse for student version and teachertrue for teachers version

% --- Solution environment ---
\newtcolorbox{solutionbox}{colback=gray!10,colframe=black!60!black,arc=1mm,boxrule=1pt}

\ifteacher
\newenvironment{solution}{%
  \begin{solutionbox}\textbf{Solution:}
}{%
  \end{solutionbox}
}
\else
  \excludecomment{solution}
\fi


\begin{document}

\begin{titlepage}

% --- Top spacing ---
\vspace*{2cm}

% --- Title page ---
\begin{center}
    {\Huge \bfseries \color{maincolor} Debate\par}
    \vspace{0.5cm}
    
    {\Large \color{accentcolor} Teaching with AI (TAI) \par}
    
    \vspace{1.5cm}
    
    \begin{center}
{\large
Mirte van der Eyden$^{1}$, Andrea López Incera$^{2}$ \\[0.3cm]

{\small
$^{1}$ Institut für Fachdidaktik, Universität Innsbruck \\ 
mirte.van-der-eyden@uibk.ac.at \\[0.2cm]


$^{2}$ Departament de didàctica de la matemàtica i les ciències experimentals, Universitat Autònoma de Barcelona \\ 
andrea.lopez.incera@uab.cat
}
}
\end{center}

\end{center}

\vfill

% --- Logos at the bottom ---
\begin{center}
\begin{tabular}{ccc}
\includegraphics[height=1.5cm]{Figures/Logos/universitaet-innsbruck-logo-cmyk-farbe.jpg} &
\includegraphics[height=1.5cm]{Figures/Logos/logotipuab-v1-verd.png} &
\includegraphics[height=1.5cm]{Figures/Logos/EN Co-funded by the EU_POS.jpg}
\end{tabular}
\end{center}


\vspace{0.5cm}

% Optional footer text
\begin{center}
    {\small Project funded by Erasmus+ Programme of the European Union}
\end{center}

\end{titlepage}

% --- Description of activity ---
\vspace*{0.5\textheight}

Subject: Debate about AI

Activity duration: 180 min

\vspace{0.5cm}

Steps:
\begin{enumerate}
    \item Distribute debate themes and prepare in groups (90 min)
    \item Debate in groups (90 min)
\end{enumerate}

Learning goals:

\begin{itemize}
    \item Doing research about an ethical topic
    \item Preparing arguments
    \item Debate
\end{itemize}

\newpage

% --- HEADER / BRANDING ---
\begin{center}
    {\Huge \color{maincolor} Debate}\\[5pt]
    {\large \textit{on ethics and practical consequences of AI}}\\
    \rule{\linewidth}{1pt}
\end{center}

\vspace{0.5cm}

Here we share a list of suggested topics for a debate about AI. Students get assigned a thesis and a side (for/against) in groups of 2 or 3. They get time to do research and prepare their arguments. 
Then, they will have a debate with the other side. The students in de audience each prepare a question. Afterwards, two students in the audience are appointed to ask a question. 


\section*{Debate Topics on AI}


Each topic has a short thesis, a simple explanation, and two sides (For / Against). Students argue for or against the thesis.


\begin{enumerate}
\item \textbf{Thesis: AI-made works should be allowed to have copyright.}\\
\emph{What this means:} If a human uses AI to make images, music, or text, can that human own the rights like an author?\\
\textbf{For:}
\begin{itemize}
\item The human still makes choices: prompts, edits, and selection show creativity.
\item It motivates people to invest time and money in creative AI.
\item Tools have always changed art (like cameras); authorship can include new tools.
\end{itemize}
\textbf{Against:}
\begin{itemize}
\item Copyright is for human authors; fully machine-made work should not qualify.
\item Could create a flood of protected machine content and legal problems.
\item Other options (contracts, licenses) may fit better than copyright.
\end{itemize}


\item \textbf{Thesis: The environmental cost of generative AI is bigger than its benefits.}\\
\emph{What this means:} Training and using big AI models uses a lot of electricity and water. Is it worth it?\\
\textbf{For:}
\begin{itemize}
    \item Data centers use much energy and water; local communities pay the price.
    \item Many AI uses are low-value or replace simpler methods.
    \item The same resources could support health, climate, or public services.
\end{itemize}
\textbf{Against:}
\begin{itemize}
    \item AI can save energy in other areas (smart grids, transport, buildings).
    \item Chips and data centers are getting more efficient.
    \item Rules and clean energy can reduce the harm without stopping progress.
\end{itemize}

\item \textbf{Thesis: AI companions and creative tools are good for mental health.}\\
\emph{What this means:} Chatbots and AI art/music tools might help people feel better.\\
\textbf{For:}
\begin{itemize}
    \item Easy to access, no stigma; good for journaling and reflection.
    \item Can reduce loneliness and support coping skills.
    \item Early studies show mood benefits when used with care.
\end{itemize}
\textbf{Against:}
\begin{itemize}
    \item People may become dependent and talk less to real friends.
    \item Bots can give wrong or unsafe advice.
    \item Privacy risks: sensitive data could be misused.
\end{itemize}

\item \textbf{Thesis: Students should be allowed to use generative AI in arts and humanities assignments.}\\
\emph{What this means:} Can students use tools like chatbots or image generators for class work?\\
\textbf{For:}
\begin{itemize}
    \item Matches real creative work today; teaches how to guide and improve AI output.
    \item Helps with brainstorming and feedback for all students.
    \item With clear rules and citation, it supports learning.
\end{itemize}
\textbf{Against:}
\begin{itemize}
    \item Makes it hard to judge a student’s own skills.
    \item Tempts students to shortcut and weakens basics (writing, analysis, technique).
    \item Not all students have equal access; tools can be biased.
\end{itemize}

\item \textbf{Thesis: Using public internet content to train AI is fair use.}\\
\emph{What this means:} Is it okay for companies to scrape online text and images to train AI without asking each creator?\\
\textbf{For:}
\begin{itemize}
    \item The training is transformative (learning patterns, not copying exact works).
    \item Needed for useful AI; similar to text and data mining.
    \item Society benefits from better tools and knowledge.
\end{itemize}
\textbf{Against:}
\begin{itemize}
    \item It takes creators’ work without permission or pay.
    \item AI outputs can replace the original creators’ market.
    \item Publicly available does not mean free to use; privacy and website rules matter.
\end{itemize}

\item \textbf{Thesis: AI-made art is art.}\\
\emph{What this means:} Should works made with or by AI be accepted as ``art''?\\
\textbf{For:}
\begin{itemize}
    \item Art depends on intention, context, and audience—not only on the tool.
    \item Humans still choose prompts, curate results, and set meaning.
    \item Art history shows new media (photo, video, code) become accepted.
\end{itemize}
\textbf{Against:}
\begin{itemize}
    \item AI lacks human feeling and agency inside the work.
    \item Often imitates styles; may lead to sameness and pastiche.
    \item Can shift focus from human craft to automation.
\end{itemize}

\item \textbf{Thesis: Teachers should use AI to make teaching materials and help evaluate students.}\\
\emph{What this means:} Is it good for teachers to use AI to prepare lessons, quizzes, and feedback?\\
\textbf{For:}
\begin{itemize}
    \item Saves time so teachers can focus on students.
    \item Helps personalize learning and improve accessibility.
    \item Can support clear rubrics and spot learning gaps.
\end{itemize}
\textbf{Against:}
\begin{itemize}
    \item AI can be biased or wrong; grading may become less fair or clear.
    \item Risks to student privacy and data security.
    \item Could weaken teacher skills and relationships with students.
\end{itemize}

\item \textbf{Thesis: In creative and cultural jobs, AI will cause more job losses than new jobs.}\\
\emph{What this means:} Will AI take away more jobs (designers, writers, translators, etc.) than it creates?\\
\textbf{For:}
\begin{itemize}
    \item Fast replacement in illustration, content writing, and post-production.
    \item Lower pay and fewer entry-level opportunities.
    \item Big tech firms capture most value, not individual creators.
\end{itemize}
\textbf{Against:}
\begin{itemize}
    \item New jobs appear (model trainers, prompt experts, curators, verifiers).
    \item Higher productivity can grow markets and demand.
    \item History shows tech changes jobs more than it removes them over time.
\end{itemize}

\item \textbf{Thesis: People must say when they used AI in academic or creative work.}\\
\emph{What this means:} Should students, writers, and artists disclose AI help in their work?\\
\textbf{For:}
\begin{itemize}
    \item Builds trust and protects academic honesty.
    \item Gives proper credit and shows the human’s real contribution.
    \item Helps with copyright, safety checks, and provenance.
\end{itemize}
\textbf{Against:}
\begin{itemize}
    \item Hard to enforce when AI is built into many tools.
    \item Might discourage experimentation with useful tools.
    \item Some say the result matters more than which tool was used.
\end{itemize}

\end{enumerate}


\subsection*{Optional tips for students}
\begin{itemize}
\item Define key words at the start in 1--2 simple sentences (e.g., copyright,'' fair use,'' ``disclosure'').
\item Use clear examples: a news article on AI energy use, a school policy, a court case, a real artist’s experience.
\item Agree on team roles: one introduces, one gives evidence, one rebuts and closes.
\end{itemize}

\end{document}